\section{Protocolo de evaluación y gestión de recursos}
\subsection{Tiempo y costo}
La preparación del instructor y el inicio del estudiante deben medirse por separado. Para cada sesión se propone registrar un identificador anónimo, hora de inicio, primera ejecución válida, resultado de la comprobación y categoría de incidencia. El informe debe incluir participantes previstos, casos medidos, ejecuciones correctas y datos ausentes, junto con mediana y dispersión de los tiempos disponibles. Una preparación previa o una imagen precargada se declara como condición de la medición.

El registro económico distingue horas activas por tamaño de máquina, horas-núcleo, almacenamiento en GB-hora y uso de Actions para preconstrucciones. El almacenamiento sigue existiendo cuando el entorno está detenido; la tarifa y las cuotas aplicables se registran con su fecha \cite{billing}. El costo por participante requiere indicar si el denominador son inscritos, asistentes o personas que completaron la práctica. El repositorio adjunta plantillas vacías y un agregador que calcula consumos a partir de registros explícitos, sin asignar una tarifa universal ni reconstruir consumos históricos ausentes.

\subsection{Fallos que el diseño debe contemplar}
El cuadro~\ref{tab:riesgos} vincula riesgos operativos con señales y respuestas. Son escenarios de diseño, no incidentes atribuidos a una cohorte. Una prueba previa debe incluir cuenta de estudiante, permisos, arranque limpio y descarga de la muestra de datos. Con una alternativa local o un cuaderno exportable se puede continuar una parte de la actividad ante interrupciones del servicio.

\begin{table}[tbp]
\caption{Riesgos de diseño y respuesta propuesta. Son escenarios de diseño, no un registro de incidentes observados.}\label{tab:riesgos}
\begin{tabularx}{\linewidth}{@{}>{\RaggedRight\arraybackslash}p{34mm}YY@{}}
\toprule
\cabecera{Riesgo} & \cabecera{Señal a registrar} & \cabecera{Respuesta prevista} \\
\midrule
\fila{Acceso o permisos} & Inicio rechazado; categoría del error & Probar cuentas y acceso antes de clase \\
\addlinespace
\fila{Cuota o presupuesto} & Aprovisionamiento bloqueado & Revisar límites y responsable de facturación \\
\addlinespace
\fila{Conectividad} & Descarga incompleta o sesión interrumpida & Muestra reducida y materiales descargables \\
\addlinespace
\fila{Dependencia incompatible} & Validación fallida tras actualizar & Imagen identificada y prueba desde estado limpio \\
\addlinespace
\fila{Pérdida de trabajo} & Entorno eliminado sin exportación & Guardado en Git y copia de resultados \\
\bottomrule
\end{tabularx}
\end{table}


La evaluación del aprendizaje necesita una tarea adicional: interpretar una figura, justificar una selección o reproducir un análisis días después. El tiempo de aprovisionamiento mide acceso al entorno; la explicación del resultado mide otra dimensión. Un diseño comparativo debería mantener tarea y condiciones equivalentes y registrar experiencia previa, conectividad y apoyo recibido.
